\documentclass[12pt]{article}
\usepackage[utf8]{inputenc}
\usepackage[T1]{fontenc}
\usepackage[margin=1in]{geometry}
\usepackage{setspace}
\usepackage{xcolor}
\usepackage[colorlinks=true, urlcolor=blue!60!black, linkcolor=blue!60!black]{hyperref}
\setlength{\parindent}{2em}
\setlength{\parskip}{1.25em}
\renewcommand{\baselinestretch}{1.0}
\begin{document}
\begin{center}
{\LARGE \textit{Korematsu v. United States}, 323 U.S. 214 (1944)\par}
\vspace{0.8em}
{\large Case Brief --- Four-Page Edition\par}
\end{center}
\vspace{1em}

\section{Memory Jogger}

Deferring to the combined wartime authority of the President and Congress, the Court upheld the conviction of an American citizen of Japanese ancestry for defying a World War II West Coast exclusion order---a decision later repudiated as ``gravely wrong the day it was decided'' (\href{https://www.courtlistener.com/opinion/104040/korematsu-v-united-states/}{323 U.S. 214}).

\section{Detailed Case Facts}

After Pearl Harbor, President Roosevelt signed Executive Order 9066 on February 19, 1942, authorizing military commanders to prescribe areas from which ``any or all persons may be excluded.'' General John L. DeWitt designated the Pacific Coast states as military areas, and on March 21, 1942, Congress made it a federal misdemeanor to remain in or violate restrictions in any such area. On March 27 DeWitt barred persons of Japanese ancestry from leaving Military Area No.~1; on May 3 Civilian Exclusion Order No.~34 excluded all persons of Japanese ancestry, alien and citizen, from part of that area, including Alameda County, after May 9, directing them to report for transfer to Assembly Centers.

Fred Korematsu, a native-born citizen and welder in San Leandro whose loyalty was never questioned, stayed in his home community past the deadline and was arrested on May 30, 1942.

\section{Procedural History}

Korematsu was convicted after a bench trial and placed on five years' probation, after which military authorities sent him to an Assembly Center and then a Relocation Center. After the Supreme Court held the probation order appealable (319 U.S. 432 (1943)), the Ninth Circuit affirmed. The Supreme Court affirmed on December 18, 1944. In 1984 Judge Marilyn Hall Patel vacated the conviction on \textit{coram nobis} after suppressed documents showed the government had misled the Court; Congress apologized and paid reparations in 1988; the Solicitor General confessed error in 2011; and \textit{Trump v. Hawaii} (2018) declared that \textit{Korematsu} ``has no place in law under the Constitution.''

\section{Judicial Votes}

Affirmed, \textbf{6--3}.
\begin{itemize}
\item \textbf{Majority:} Justice Black, joined by Chief Justice Stone and Justices Reed, Frankfurter, Douglas, and Rutledge.
\item \textbf{Concurrence:} Justice Frankfurter.
\item \textbf{Dissents:} Justices Roberts, Murphy, and Jackson, each writing separately.
\end{itemize}

\section{Holding}

Legal restrictions curtailing the civil rights of a single racial group are ``immediately suspect'' and subject to ``the most rigid scrutiny,'' justifiable only by ``pressing public necessity.'' Yet when Congress and the President act jointly under the war powers and military authorities apprehend ``the gravest imminent danger to the public safety,'' temporary exclusion of citizens of Japanese ancestry from a threatened area is permissible. The holding covered only the exclusion order as of May 1942, not detention in the camps, addressed the same day in \textit{Ex parte Endo}. Ironically, the ``immediately suspect'' language became the foundation of modern strict scrutiny though the Court applied no meaningful scrutiny.

\section{Analysis of Opinions}

\textbf{Black (Court).} Relying on \textit{Hirabayashi}, Black would not ``reject as unfounded the judgment of the military authorities and of Congress'' that disloyal persons could not be quickly separated from the loyal. The exclusion order superseded the earlier order forbidding departure, and because Korematsu was not convicted of refusing detention, the Court passed only on exclusion. Black denied that exclusion stemmed from racial hostility and refused to second-guess the military ``by availing ourselves of the calm perspective of hindsight.''

\textbf{Frankfurter (concurring).} The war power is ``the power to wage war successfully''; a valid military order may be enforced in civil courts, and finding it constitutional does not mean approving it.

\textbf{Roberts (dissenting).} Korematsu faced ``two diametrically contradictory orders,'' a ``cleverly devised trap'' to imprison him based solely on ancestry, violating due process.

\textbf{Murphy (dissenting).} Courts must ask whether a deprivation is reasonably related to an ``immediate, imminent, and impending'' danger. DeWitt's Final Report rested on racial stereotypes; no person of Japanese ancestry had been charged with espionage or sabotage, months passed before removal, and individual loyalty hearings were feasible.

\textbf{Jackson (dissenting).} Korematsu was punished for presence in his home state because of ancestry, though ``guilt is personal and not inheritable.'' Courts cannot assess military necessity, but they must not ``distort the Constitution''; judicial validation of the order creates a lasting principle, so civil courts should refuse to enforce it.

\section{Examples: Future Applications}

\begin{enumerate}
\item \textbf{Same-Side (Invalid): Group Detention of Citizens.} An order requiring citizens of one national origin near coastal bases to report to holding centers without hearings is ``objectively unlawful and outside the scope of Presidential authority'' under \textit{Trump v. Hawaii} and \textit{Endo}.
\item \textbf{Same-Side (Invalid): Detention on Bare Assertion.} Holding a citizen without charge on a commander's claim of ``pressing public necessity'' fails under \textit{Hamdi v. Rumsfeld}.
\item \textbf{Opposite-Side (Valid): Foreign-Entry Restrictions.} A statutorily authorized suspension of entry for nationals of certain countries was upheld in \textit{Trump v. Hawaii}, which distinguished \textit{Korematsu}.
\item \textbf{Opposite-Side (Valid): Neutral Emergency Evacuation.} A 48-hour evacuation of all residents from a blast zone is permissible; even Roberts accepted temporary, non-discriminatory exclusions.
\item \textbf{Fence-Sitter: Travel-Based Screening.} Heightened screening of citizens who visited a hostile nation, falling heavily on one ethnic community, pits group-suspicion concerns against deference to neutral, congressionally authorized criteria.
\end{enumerate}

\section{Critique}

\textbf{Originalist and textualist views.} Scholars place \textit{Korematsu} in the ``anticanon'' with \textit{Dred Scott} and \textit{Plessy}, arguing that due process, the jury right, and the Suspension Clause forbid detaining citizens outside the criminal process, and faulting the standardless delegation to military commanders.

\textbf{Progressive views.} Civil-rights scholars stress the racial animus behind Executive Order 9066 and the failure of judicial review when officials withhold evidence. Harold Hongju Koh argues that \textit{Trump v. Hawaii}, while repudiating \textit{Korematsu}, risked resurrecting its ``ghost'' through excessive deference.

\section{Key Quotations}

\begin{quote}
``It should be noted, to begin with, that all legal restrictions which curtail the civil rights of a single racial group are immediately suspect. That is not to say that all such restrictions are unconstitutional. It is to say that courts must subject them to the most rigid scrutiny. Pressing public necessity may sometimes justify the existence of such restrictions; racial antagonism never can.''
\end{quote}
\noindent --- Justice Black, Opinion of the Court.

\begin{quote}
``But once a judicial opinion rationalizes such an order to show that it conforms to the Constitution, or rather rationalizes the Constitution to show that the Constitution sanctions such an order, the Court for all time has validated the principle of racial discrimination in criminal procedure and of transplanting American citizens. The principle then lies about like a loaded weapon ready for the hand of any authority that can bring forward a plausible claim of an urgent need.''
\end{quote}
\noindent --- Justice Jackson, dissenting.

\begin{quote}
``This exclusion of `all persons of Japanese ancestry, both alien and non-alien,' from the Pacific Coast area on a plea of military necessity in the absence of martial law ought not to be approved. Such exclusion goes over `the very brink of constitutional power' and falls into the ugly abyss of racism.''
\end{quote}
\noindent --- Justice Murphy, dissenting.

\end{document}
